\documentclass[10pt,aspectratio=169]{beamer}
\input{../../_en_preambulo_beamer}
% Comandos y macros estadísticas compartidas para las presentaciones Beamer en inglés (Metropolis)

% Conjuntos numéricos básicos
\newcommand{\R}{\mathbb{R}}
\newcommand{\N}{\mathbb{N}}
\newcommand{\Q}{\mathbb{Q}}
\newcommand{\Z}{\mathbb{Z}}
\newcommand{\set}[1]{\left\{ #1 \right\}}
\newcommand{\abs}[1]{\left|#1\right|}
\newcommand{\norm}[1]{\left\| #1 \right\|}

% Comandos estadísticos y probabilísticos del curso
\newcommand{\Var}{\operatorname{Var}}
\newcommand{\cov}{\operatorname{Cov}}
\newcommand{\corr}{\rho}
\newcommand{\comb}[2]{\begin{pmatrix} #1 \\ #2 \end{pmatrix}}
\newcommand{\s}{\sigma}
\newcommand{\particion}{\mathfrak{P}}
\newcommand{\card}[1]{\#\left( #1 \right)}
\newcommand{\seq}[3]{\set{#1}_{#2}^{#3}}
\newcommand{\E}{\mathbb{E}}
\newcommand{\Prob}{\mathbb{P}}

% Entornos pedagógicos y cajas en inglés académico natural (soportando tanto nombres en inglés como en español)
\newenvironment{discussion}{%
  \begin{alertblock}{Discussion Question}%
}{%
  \end{alertblock}%
}
\newenvironment{discusion}{%
  \begin{alertblock}{Discussion Question}%
}{%
  \end{alertblock}%
}

\newenvironment{reflection}{%
  \begin{alertblock}{Classroom Reflection}%
}{%
  \end{alertblock}%
}
\newenvironment{reflexion}{%
  \begin{alertblock}{Classroom Reflection}%
}{%
  \end{alertblock}%
}

\newenvironment{paradox}{%
  \begin{alertblock}{Exploratory Paradox}%
}{%
  \end{alertblock}%
}
\newenvironment{paradoja}{%
  \begin{alertblock}{Exploratory Paradox}%
}{%
  \end{alertblock}%
}

\newenvironment{motivation}{%
  \begin{exampleblock}{Motivation: Data Science in Practice}%
}{%
  \end{exampleblock}%
}
\newenvironment{motivacion}{%
  \begin{exampleblock}{Motivation: Data Science in Practice}%
}{%
  \end{exampleblock}%
}

\newenvironment{quickchallenge}{%
  \begin{block}{Quick Team Challenge}%
}{%
  \end{block}%
}
\newenvironment{retoclase}{%
  \begin{block}{Quick Team Challenge}%
}{%
  \end{block}%
}


\title[Proportion Intervals]{Section 6.5: Estimation of a Proportion and Difference Between Two Proportions}
\subtitle{Wald, Wilson, and rate comparisons}
\author[J. Castillo Colmenares]{Juliho Castillo Colmenares}
\institute{Tecnológico de Monterrey}
\date{}

\begin{document}
\begin{frame}[plain]\vspace{-0.8cm}\titlepage\end{frame}

\begin{frame}{Roadmap --- Chapter 6}
  \footnotesize
  \begin{itemize}\itemsep=0.08cm
    \item \textbf{06.02} Mean intervals
    \item \textbf{06.04} Standard errors
    \item \textbf{06.05} Proportions and differences \emph{(today)}
    \item \textbf{06.07} Sample size
  \end{itemize}
  \begin{block}{Learning objective}
    Build intervals for one proportion and for the difference of two proportions,
    distinguishing Wald from Wilson.
  \end{block}
\end{frame}

\begin{frame}{Motivation: categorical responses}
  \footnotesize
  Proportions describe Bernoulli successes: click-through rates, preferences,
  and classification metrics.
  \begin{alertblock}{Guiding question}
    How do we quantify uncertainty in $\hat p=X/n$ and in a rate difference?
  \end{alertblock}
\end{frame}

\begin{frame}{Wald interval}
  \footnotesize
  When large-sample conditions hold,
  \[
    \hat p\pm Z_{\alpha/2}\sqrt{\frac{\hat p(1-\hat p)}{n}}.
  \]
  It uses a normal approximation to the sampling distribution of $\hat p$.
\end{frame}

\begin{frame}{Wald limitation}
  \footnotesize
  For small samples or proportions near $0$ or $1$, Wald may have poor actual
  coverage and produce limits outside $[0,1]$.
  \begin{block}{Reference conditions}
    The reading uses $n\hat p\ge5$ and $n(1-\hat p)\ge5$ (or conservatively
    $\ge10$).
  \end{block}
\end{frame}

\begin{frame}{Wilson interval}
  \footnotesize
  Directly inverting the score test gives
  \[
  \frac{1}{1+Z^2/n}\left[\hat p+\frac{Z^2}{2n}\pm
  Z\sqrt{\frac{\hat p(1-\hat p)}{n}+\frac{Z^2}{4n^2}}\right].
  \]
  The interval need not be centered at $\hat p$.
\end{frame}

\begin{frame}{Difference of two proportions}
  \footnotesize
  For sufficiently large independent samples,
  \[
    (\hat p_1-\hat p_2)\pm Z_{\alpha/2}
    \sqrt{\frac{\hat p_1(1-\hat p_1)}{n_1}+
    \frac{\hat p_2(1-\hat p_2)}{n_2}}.
  \]
  The central estimator is a difference and the variance combines both groups.
\end{frame}

\begin{frame}{Conditions and method choice}
  \footnotesize
  Wald is a large-sample approximation; Wilson is preferable with small counts
  or a proportion near a boundary.
  \begin{alertblock}{Rule}
    Check conditions first, then choose the formula.
  \end{alertblock}
\end{frame}

\begin{frame}{Computational bridge 1/4: estimator}
  \footnotesize
  \[
    \hat p=\frac{X}{n},\qquad
    \operatorname{Var}(\hat p)\approx\frac{p(1-p)}{n}.
  \]
  Wald estimates the standard error by replacing $p$ with $\hat p$.
\end{frame}

\begin{frame}{Computational bridge 2/4: Wilson correction}
  \footnotesize
  Wilson keeps $p$ in the score denominator while inverting the inequality.
  \begin{block}{Consequence}
    Center and width adjust automatically near the boundaries.
  \end{block}
\end{frame}

\begin{frame}{Computational bridge 3/4: two groups}
  \footnotesize
  \[
    \widehat{SE}(\hat p_1-\hat p_2)=
    \sqrt{\frac{\hat p_1(1-\hat p_1)}{n_1}+
    \frac{\hat p_2(1-\hat p_2)}{n_2}}.
  \]
  Each sample contributes its own uncertainty.
\end{frame}

\begin{frame}{Computational bridge 4/4: boundaries}
  \footnotesize
  \begin{itemize}\itemsep=0.12cm
    \item Wald may exceed $[0,1]$.
    \item Wilson maintains plausible limits.
    \item A difference $p_1-p_2$ lies in $[-1,1]$.
  \end{itemize}
\end{frame}

\begin{frame}{Reading application 1/4 --- Wald}
  \footnotesize
  The reading gives the normal CI for $p$ when large-sample conditions hold.
  \begin{block}{Statement}
    Identify $X$, $n$, $\hat p$, and the critical value.
  \end{block}
\end{frame}

\begin{frame}{Reading application 1/4 --- solution}
  \footnotesize
  \[
    CI_W=\hat p\pm Z_{\alpha/2}
    \sqrt{\frac{\hat p(1-\hat p)}{n}}.
  \]
  The approximation is reasonable when expected successes and failures are
  sufficiently numerous.
\end{frame}

\begin{frame}{Reading application 2/4 --- Wilson}
  \footnotesize
  For a small sample or an extreme proportion, the reading recommends score
  inversion rather than Wald.
  \begin{block}{Statement}
    Replace the symmetric formula with Wilson's expression.
  \end{block}
\end{frame}

\begin{frame}{Reading application 2/4 --- solution}
  \footnotesize
  The factor $1/(1+Z^2/n)$ and the term $Z^2/(4n^2)$ correct center and width.
  \begin{alertblock}{Result}
    Wilson keeps actual coverage closer to the nominal level.
  \end{alertblock}
\end{frame}

\begin{frame}{Reading application 3/4 --- two proportions}
  \footnotesize
  Two independent populations have rates $p_1,p_2$ and sample sizes $n_1,n_2$.
  \begin{block}{Statement}
    Estimate $p_1-p_2$ with $\hat p_1-\hat p_2$.
  \end{block}
\end{frame}

\begin{frame}{Reading application 3/4 --- solution}
  \footnotesize
  \[
    CI=(\hat p_1-\hat p_2)\pm Z_{\alpha/2}\widehat{SE}.
  \]
  The interval includes zero when the observed difference is compatible with no
  difference.
\end{frame}

\begin{frame}{Reading application 4/4 --- conditions}
  \footnotesize
  Normal approximation validity depends on expected counts
  $n_i\hat p_i$ and $n_i(1-\hat p_i)$.
  \begin{block}{Statement}
    Check conditions before reporting an interval.
  \end{block}
\end{frame}

\begin{frame}{Reading application 4/4 --- decision}
  \footnotesize
  If conditions fail, prefer Wilson (or Agresti--Coull) for one proportion; for
  two rates, choose an approximation appropriate to the design.
  \begin{exampleblock}{Result}
    The method is selected from the available information, not convenience.
  \end{exampleblock}
\end{frame}

\begin{frame}{Synthesis of the section}
  \footnotesize
  \begin{itemize}\itemsep=0.12cm
    \item Wald is simple and needs large-sample conditions.
    \item Wilson corrects center, width, and coverage.
    \item Two proportions combine two standard errors.
    \item Expected counts determine the appropriate approximation.
  \end{itemize}
\end{frame}

\begin{frame}{Curricular transition}
  \footnotesize
  \begin{block}{Established}
    Proportion intervals must respect the support $[0,1]$ and the assumptions of
    the chosen approximation.
  \end{block}
  \begin{alertblock}{Next section}
    \textbf{06.07 Sample Size}: choose $n$ before observing data while controlling
    confidence and margin of error.
  \end{alertblock}
\end{frame}
\end{document}
